\documentclass[../../main.tex]{subfiles}
\begin{document}

{\bf [解答]}

コーシーシュワルツの不等式から
\begin{align}
  \sqrt{x}+\sqrt{y}
  = \dfrac{\sqrt{2x}}{\sqrt{2}}+\sqrt{y}
  \le \sqrt{\left(\dfrac{1}{2}+1\right)(2x+y)}
  = \sqrt{\dfrac{3}{2}}\sqrt{2x+y}
\end{align}
である．等号成立は
$
\begin{pmatrix}\sqrt{x}\\\sqrt{y}
\end{pmatrix}\parallel
\begin{pmatrix}\dfrac{1}{\sqrt{2}}\\1
\end{pmatrix}$
，つまり$x=\dfrac{y}{4}$の時である．
よって求める最小の$k$は
\begin{align}
  k = \sqrt{\dfrac{3}{2}} = \dfrac{\sqrt{6}}{2}
\end{align}
である．

\newpage
\end{document}
